% Drawing for measurement `mos-rdson`
% Stands for the deck lines:
%   BR0 0 dr0 I=exp(2.302585092994046*V(x))
%   DCL0 dr0 cl0 DCLMP / VCL0 cl0 0 30
%   VGR0 gr0 0 10
%   MR0 dr0 gr0 0 0 DUT
%   .dc VX -3 2 0.05
\documentclass[border=8pt]{standalone}
\usepackage[siunitx,RPvoltages,european]{circuitikz}
\ctikzset{bipoles/length=0.9cm, label/align=straight, inductors/scale=0.8}
\ctikzset{inductor=american}
\usepackage{amsmath}
\begin{document}
\begin{circuitikz}[font=\footnotesize, line width=0.5pt]
\draw (0,0) node[nmos](Q){};
\node[anchor=west] at ($(Q.D)!0.5!(Q.S) + (0.55,0)$) {DUT};
\draw (Q.S) to[short] ++(0,-0.5) node[ground]{};
\draw (Q.D) to[short] ++(0,0.3) coordinate(Qd) to[isource, l={\shortstack[l]{$I_D=10^{V(x)}$\\ \SI{1}{mA}\dots\SI{100}{A}}}] ++(0,1.8) node[ground, rotate=180]{};
\draw (Qd) to[short] ++(1.6,0) to[D, l={clamp}] ++(0,1.4) to[short] ++(0,0.2) node[vcc]{\SI{30}{V}};
\draw (Qd) to[short] ++(3.4,0) to[rmeter, t=V, l={$V_{DS}$}] ++(0,-2.0) node[ground]{};
\draw (Q.G) to[short] ++(-2.2,0) to[vsource, l_={\shortstack[r]{$V_{GS}$ step\\ 2.5\dots\SI{15}{V}}}, invert] ++(0,-1.9) node[ground]{};
\node[anchor=west, align=left] at (-4.6,-3.6) {the drain current is forced by the behavioural dial;\\ the clamp catches a copy that cannot pass it; $R_{DS(\mathrm{on})}=V_{DS}/I_D$\\ second deck: $I_D$ held at 0.1, 1, \SI{10}{A}, gate swept $0\to\SI{15}{V}$};
\end{circuitikz}
\end{document}
